\section{Fear of preventable harm: the errors we can no longer accept}
\label{sec:fear}

The second appeal is the fear of harm that did not have to happen. Missed
diagnoses, medication errors, unsafe devices, and delayed treatment have driven
drug-safety regulation, device oversight, and electronic health record
requirements for a generation. The underlying message is simple: without this
law, more people will suffer than need to.

[DRAFTING INSTRUCTIONS: in the full-narrative, anchor the fear appeal in the two
strongest credible sources in \texttt{review/references/narrative\_refs.bib}:
\texttt{kohn2000toerr} (To Err Is Human) and \texttt{makary2016medical} (medical
error as a leading cause of death). Contrast the legacy chain of manual handoffs
with the verified path using
\texttt{cancer-automated/papers/VVUQ-02} (each gate catches a defect where it
occurs). Reproduce Mermaid figures
\texttt{review/mermaid/diagrams/04-compounding-human-error.md} and
\texttt{20-patient-throughput-flow.md} as colored TikZ. Use the KEY TERM
\texttt{markup} to note where committee debate would harden the safety language.]

\begin{narrativefig}
\node[nfrisk] (a) at (0,0) {Manual handoff};
\node[nfrisk] (b) at (3.4,0) {Error compounds};
\node[nfdec] (c) at (7.0,0) {Gate};
\node[nfhope] (d) at (10.2,0) {Contained};
\draw[nfedge] (a)--(b); \draw[nfedge] (b)--(c); \draw[nfedge] (c)--(d);
\end{narrativefig}
\vspace{-0.2cm}
\figcaption{Figure 3. Compounding error versus a contained defect (placeholder;
expanded from Mermaid 04 in the full-narrative).}

A bill cannot promise that nothing will ever go wrong. It can promise that the
system checks itself at every step, so the same mistake is not made twice and is
not made invisibly.
